\section{Four-term arithmetic progressions}
\label{sec:ap4}
Let \(F=\{0,1,2,3\}\). For \(k=(k_0,k_1,k_2,k_3)\ge0\), write

\[
T(k)=T_F(k)=\sum_{s=0}^{12}
\max_{i_1+i_2+i_3+i_4=s}\prod_{r=1}^{4}k_{i_r},
\qquad i_r\in F.
\]

For a tuple \(v\), let \(v^R_i=v_{3-i}\). Define the five squared profiles

\[
\begin{aligned}
A&=(1,1/4,0,1/4),\\
B&=(1,0,1/4,1/4),\\
C&=(1,1/4,9/64,25/256),\\
D&=(1/4,1,0,1/4),\\
E&=(0,1,1/4,9/64),
\end{aligned}
\]

and the ten-element set

\[
\mathcal V=\{A,B,C,D,E,A^R,B^R,C^R,D^R,E^R\}.
\]

\begin{proposition}\label{prop:ap4}
The weighted norm has the finite-profile formula

\begin{equation}\label{eq:a1}
\boxed{\quad
\sup_{x\ge0,\;x*x\le1}\sum_{i=0}^{3}k_ix_i^2
=\max_{v\in\mathcal V}k\cdot v,
\quad}
\end{equation}

where the convolution inequality is coefficientwise at all seven indices, and

\begin{equation}\label{eq:a2}
\boxed{\quad
(k\cdot v)^4\le T(k)\qquad(v\in\mathcal V).
\quad}
\end{equation}

Consequently,

\begin{equation}\label{eq:a3}
x*x\le1\quad\Longrightarrow\quad
\left(\sum_i k_ix_i^2\right)^4\le T(k)
\qquad(k\ge0).
\end{equation}

For positive \(k_i\), comparison \eqref{eq:a2}, and hence \eqref{eq:a3} at the maximizing value,
is strict. These statements hold on every four-term integer arithmetic
progression after relabeling the coordinates.
\end{proposition}

\subsection{Feasibility and a six-profile convex hull}

Each member of \(\mathcal V\) is the square of a feasible tuple. For the five
unreversed members, the seven coefficients of \(\sqrt v*\sqrt v\) are

\begin{center}
\begin{tabular}{ll}
\toprule
Profile & Ordinary convolution coefficients\\
\midrule
\(A\) & \((1,1,1/4,1,1/2,0,1/4)\)\\
\(B\) & \((1,0,1,1,1/4,1/2,1/4)\)\\
\(C\) & \((1,1,1,1,29/64,15/64,25/256)\)\\
\(D\) & \((1/4,1,1,1/2,1,0,1/4)\)\\
\(E\) & \((0,0,1,1,1,3/8,9/64)\)\\
\bottomrule
\end{tabular}
\end{center}

Reversal reverses this coefficient list. Thus the right-hand side of \eqref{eq:a1} is
always attained by a feasible squared profile.

Let

\[
\mathcal S=\{A,B,D,A^R,B^R,D^R\}.
\]

Its convex hull has the exact elementary description

\begin{equation}\label{eq:a4}
\operatorname{conv}(\mathcal S)=
\left\{y\ge0:
\sum_i y_i=\frac32,\quad y_0,y_3\ge\frac14,\quad
y_0+y_3\le\frac54\right\}.
\end{equation}

To check the reverse inclusion, put

\[
r_0=\frac{4y_0-1}{3},\qquad
r_3=\frac{4y_3-1}{3},\qquad r_m=1-r_0-r_3.
\]

These are nonnegative and sum to one. Allocate the same fraction
\(p=y_1/(y_1+y_2)\) to the first member of each pair
\((A,B)\), \((B^R,A^R)\), and \((D,D^R)\), with pair masses
\(r_0,r_3,r_m\), respectively. The resulting convex combination is \(y\).
The denominator is positive because \(y_1+y_2\ge1/4\).

A useful consequence is the following sufficient domination test. If a tuple
\(b\ge0\) satisfies

\[
b_0,b_3\le\frac12,\qquad m=b_1+b_2\le1,
\]

\begin{equation}\label{eq:a5}
m+\max(b_0,1/4)+\max(b_3,1/4)\le\frac32,
\end{equation}

then some \(y\in\operatorname{conv}(\mathcal S)\) dominates \(b\)
coordinatewise. Indeed choose the final middle sum to be
\(m'=\max(m,1/4)\), and the endpoint sum to be \(3/2-m'\). The latter lies
in \([1/2,5/4]\) and is at least
\(\max(b_0,1/4)+\max(b_3,1/4)\): this is \eqref{eq:a5} when \(m\ge1/4\), and
follows from \(b_0,b_3\le1/2\) otherwise. Increase coordinates to these sums
and apply \eqref{eq:a4}.

\subsubsection{Self-contained proof of the sharp norm lemma}

Let \(N=\sum_i x_i^2\) and \(M=\max_i x_i^2\). The constraints imply
\(x_i^2\le1\), \(2x_ix_j\le1\) for \(i\ne j\), and
\(x_0x_3+x_1x_2\le1/2\). We repeatedly use

\begin{equation}\label{eq:a6a}
0\le r,s\le J\quad\Longrightarrow\quad r+s\le J+rs/J,
\end{equation}

which follows from \((J-r)(J-s)\ge0\).

If \(M\le1/2\), put \(r=x_0x_3\), \(s=x_1x_2\). Applying \eqref{eq:a6a} to the
squared coordinates in each pair yields

\[
N\le1+2r^2+2s^2\le1+2(r+s)^2\le3/2.
\]

Suppose \(1/2\le M\le3/4\), attained at index \(i\). Put
\(y=x_ix_{3-i}\), \(J=1/(4M)\), and let \(r\) be the product of the two
coordinates outside this pair. The other three squared coordinates are at most
\(J\), while \(r+y\le1/2\) and \(r\le J\). Hence

\[
N\le M+y^2/M+J+4M\min(J,1/2-y)^2.
\]

For \(0\le y\le1/2-J\), this expression increases with \(y\). For
\(1/2-J\le y\le1/2\), it is a convex quadratic. Its maximum is therefore
the larger of the two endpoint values

\[
B(M)=M+\frac1{2M}+\frac{(2M-1)^2}{16M^3},\qquad M+\frac1{2M},
\]

namely \(B(M)\). With \(w=2M-1\in[0,1/2]\),

\[
16M^3(3/2-B(M))=w(1-w^2-w^3)\ge0.
\]

Suppose \(M\ge3/4\) is at index one; index two follows by reversal. Put
\(J=1/(4M)\). The constraint \(M+2x_0x_2\le1\) and \eqref{eq:a6a} imply

\[
x_0^2+x_2^2\le J+M(1-M)^2,
\qquad N\le M+\frac1{2M}+M(1-M)^2.
\]

The difference from \(3/2\) is

\[
\frac{(1-M)(2M^3-2M^2+2M-1)}{2M}\ge0.
\]

The cubic increases, since its derivative is
\(6(M-1/3)^2+4/3>0\), and is \(7/32\) at \(M=3/4\).

Finally suppose an endpoint has squared value \(u\in[1/2,1]\). By reversal
take \(x_0^2=u\), and put
\(a=2x_0x_1\), \(b=2x_0x_2\), \(c=2x_0x_3\). The first four constraints
give

\[
0\le a\le1,\qquad 0\le b\le1-a^2/4,\qquad 0\le c\le1-ab/2.
\]

For fixed \(a\), the upper bound \(a^2+b^2+(1-ab/2)^2\) on
\(a^2+b^2+c^2\) is convex in \(b\). At \(b=0\), it is at most two. At
\(b=1-a^2/4\), it is

\[
\Psi(a)=2-a+3a^2/4+a^3/4-a^4/16+a^6/64.
\]

This is convex on \([0,1]\), since
\(\Psi''(a)=3/2+3a/2-3a^2/4+15a^4/32\ge3/4\), and its endpoint values
are \(2\) and \(125/64<2\). Therefore

\[
N=u+\frac{a^2+b^2+c^2}{4u}\le u+\frac1{2u}\le\frac32.
\]

These cases cover every feasible tuple. The norm bound is sharp at
\(x=(1/2,1,0,1/2)=\sqrt D\), as its feasibility table shows.

\subsection{Reduction to the ten profiles}

For a feasible \(x\), write \(b_i=x_i^2\). The convolution constraints give

\begin{equation}\label{eq:a6}
b_i\le1,\qquad b_ib_j\le\frac14\quad(i\ne j).
\end{equation}

We also use the previously proved sharp norm lemma

\begin{equation}\label{eq:a7}
\sum_i b_i\le\frac32.
\end{equation}

\subsubsection{All squared coordinates at most one half}

Suppose \(b_i\le1/2\) for all \(i\). Each consecutive triple has sum at most
\(9/8\). For example, the constraint
\(b_1+2\sqrt{b_0b_2}\le1\), together with \(b_0,b_2\le1/2\), gives

\[
b_0+b_2\le\frac12+2b_0b_2
\le\frac12+\frac{(1-b_1)^2}{2},
\]

\begin{equation}\label{eq:a8}
b_0+b_1+b_2\le1+\frac{b_1^2}{2}\le\frac98.
\end{equation}

The other triple is the reversed argument.

Condition \eqref{eq:a5} now holds. If both endpoints exceed \(1/4\), use \eqref{eq:a7}; if both
are below \(1/4\), use \(b_1+b_2\le1\). If just one is below \(1/4\),
the expression in \eqref{eq:a5} is a consecutive triple sum plus \(1/4\), at most
\(11/8<3/2\). Thus \(b\) is dominated by the convex hull of \(\mathcal S\).

\subsubsection{A middle squared coordinate at least one half}

By reversal, it suffices to write

\[
b=(a,u,v,d),\qquad 1/2\le u\le1.
\]

Put \(J=1/(4u)\). By \eqref{eq:a6}, \(a,v,d\le J\le1/2\). The constraint at
index two gives \(\sqrt{av}\le(1-u)/2\). Consequently

\[
a+v\le J+\frac{av}{J}\le\frac1{4u}+u(1-u)^2.
\]

The following bound is useful throughout this case:

\begin{equation}\label{eq:a9}
a+u+v\le\frac54.
\end{equation}

Indeed
\[
\frac54-\left(u+\frac1{4u}+u(1-u)^2\right)
=\frac{(1-u)(4u^3-4u^2+4u-1)}{4u}\ge0.
\]

The cubic is increasing, because \(12u^2-8u+4>0\), and its value at
\(u=1/2\) is \(1/2\).

First suppose \(m=u+v\le1\). The domination test \eqref{eq:a5} holds. The cases when
both endpoints are on the same side of \(1/4\) follow from \eqref{eq:a7} or \(m\le1\).
If \(a\ge1/4\) and \(d\le1/4\), use \eqref{eq:a9}. If \(a\le1/4\) and
\(d\ge1/4\), the index-four constraint gives

\[
u+v+d\le1+(\sqrt u-\sqrt d)^2\le\frac54,
\]

because \(d\le1/2\le u\le1\) and \(d\ge1/4\), so
\(0\le\sqrt u-\sqrt d\le1/2\). Thus \eqref{eq:a5} again proves domination by
\(\operatorname{conv}(\mathcal S)\).

Now suppose \(m=u+v>1\). Define

\[
\theta=4(m-1),\qquad \beta=4-3u-4v,\qquad \gamma=1-u.
\]

All three numbers are nonnegative and their sum is one. For \(\beta\), use
\(v\le1/(4u)\) to obtain

\[
\beta\ge4-3u-1/u=\frac{(1-u)(3u-1)}{u}\ge0.
\]

The convex combination

\begin{equation}\label{eq:a10}
y=\theta E+\beta D+\gamma D^R
=\left(\frac54-m,\ u,\ v,\ \frac{11-7m}{16}\right)
\end{equation}

dominates \(b\). The first coordinate does so by \eqref{eq:a9}. For the last, the
index-four constraint gives \(d\le(1-v)^2/(4u)\). We claim

\begin{equation}\label{eq:a11}
\frac{(1-v)^2}{4u}\le\frac{11-7u-7v}{16}
\qquad\left(1-u\le v\le\frac1{4u}\right).
\end{equation}

After multiplying by \(16u\), the required nonnegative expression is

\[
G_u(v)=11u-7u^2-7uv-4+8v-4v^2.
\]

It is concave in \(v\), so its minimum on this interval is at an endpoint.
The two endpoint values are

\[
G_u(1-u)=4u(1-u),
\]

\[
G_u(1/(4u))
=\frac{(1-u)(28u^3-16u^2+7u-1)}{4u^2}.
\]

The latter cubic is increasing on \([1/2,1]\), since its derivative
\(84u^2-32u+7\) is positive there, and its value at \(1/2\) is two.
This proves \eqref{eq:a11}, then \eqref{eq:a10}, and completes the middle-coordinate case.

\subsubsection{An endpoint squared coordinate at least one half}

By reversal, assume \(u=x_0^2\in[1/2,1]\). Put

\[
\alpha=2x_0x_1,\qquad \beta=2x_0x_2,\qquad \gamma=2x_0x_3.
\]

The first four convolution constraints give

\[
0\le\alpha\le1,\qquad
0\le\beta\le1-\alpha^2/4,\qquad
0\le\gamma\le1-\alpha\beta/2.
\]

For arbitrary nonnegative \(k\),

\[
k\cdot b
=k_0u+\frac{k_1\alpha^2+k_2\beta^2+k_3\gamma^2}{4u}.
\]

For fixed \(\alpha\), the numerator is at most

\[
k_1\alpha^2+k_2\beta^2+k_3(1-\alpha\beta/2)^2,
\]

which is convex in \(\beta\). At \(\beta=0\), it is at most \(k_1+k_3\).
At \(\beta=1-\alpha^2/4\), it is the polynomial

\[
\begin{aligned}
\Phi(\alpha)={}&k_2+k_3-k_3\alpha
+(k_1-k_2/2+k_3/4)\alpha^2
+k_3\alpha^3/4\\
&+(k_2/16-k_3/8)\alpha^4+k_3\alpha^6/64.
\end{aligned}
\]

This polynomial has its maximum on \([0,1]\) at an endpoint, even when it is
not convex. To see this, observe

\[
\Phi'(0)=-k_3\le0,
\]

\[
\Phi'''(\alpha)=\frac32 k_2\alpha
+\frac38 k_3(4-8\alpha+5\alpha^3)\ge0.
\]

For the last inequality, the minimum of the cubic on \([0,1]\) is at
\(\alpha=\sqrt{8/15}\), where its value is
\(4-(16/3)\sqrt{8/15}\ge0\), since \(\sqrt{8/15}\le3/4\).
Thus \(\Phi'\) is convex and starts nonpositive; its nonpositive sublevel
set is an interval containing zero. Accordingly \(\Phi\) first decreases
and then increases, with either portion possibly absent. Its maximum is at an
endpoint.

The two endpoint values of \(\Phi\) are \(k_2+k_3\) and
\(k_1+9k_2/16+25k_3/64\). Therefore the numerator is at most

\[
H=\max\{k_1+k_3,\ k_2+k_3,\ k_1+9k_2/16+25k_3/64\}.
\]

For each of these three linear forms, convexity in \(u\in[1/2,1]\) reduces
\(k_0u+H/(4u)\) to its two endpoint values. At \(u=1\), the resulting
squared profiles are \(A,B,C\). At \(u=1/2\), they are

\[
W_A=(1/2,1/2,0,1/2),\qquad
W_B=(1/2,0,1/2,1/2),
\]

\[
W_C=(1/2,1/2,9/32,25/128).
\]

The first two belong to the sparse convex hull, since

\[
W_A=(A+D+B^R)/3,\qquad
W_B=(B+D^R+A^R)/3.
\]

The last is dominated coordinatewise by

\[
\frac13 C+\frac5{12}D+\frac14D^R
=\left(\frac12,\frac12,\frac{19}{64},\frac{51}{256}\right)
\ge W_C.
\]

Consequently \(k\cdot b\le\max_{v\in\mathcal V}k\cdot v\) in this case
as well. Every feasible tuple falls into one of the three cases. Together with
the explicit feasibility table, this proves \eqref{eq:a1}.

\subsection{Comparing the ten linear forms with the four-fold maximum}

Reversal preserves \(T\), so it suffices to handle \(A,B,C,D,E\).

\subsubsection{A coefficient lemma; profiles C and E}

If \(v\ge0\) has the property that every coefficient of
\((\sum_i v_i z^i)^4\) is at most one, then

\begin{equation}\label{eq:a12}
(k\cdot v)^4
=\sum_s\sum_{i_1+\cdots+i_4=s}\prod_r(v_{i_r}k_{i_r})
\le\sum_s\max_{i_1+\cdots+i_4=s}\prod_r k_{i_r}=T(k).
\end{equation}

The nonnegative formal series

\[
(1-z)^{-1/4}=1+z/4+5z^2/32+15z^3/128+\cdots
\]

coefficientwise dominates
\(1+z/4+9z^2/64+25z^3/256\). Raising to the fourth power preserves this
order, and \(((1-z)^{-1/4})^4=(1-z)^{-1}\) has all coefficients equal to
one. This proves the required coefficient property for \(C\).

For \(E\), use
\(\sum_i E_i z^i=z(1+z/4+9z^2/64)\). The same comparison applies before
the harmless fourth-power shift by four. Thus \eqref{eq:a12} proves both profiles.

\subsubsection{Profile B}

Set \(a=k_0\), \(b=k_2\), \(c=k_3\), and ignore \(k_1\); doing so can
only decrease \(T\). The target linear form is \(a+b/4+c/4\), on indices
\(0,2,3\). Grouping its fourth-power expansion by index sum, the only fibers
with more than one monomial are

\begin{center}
\begin{tabular}{lll}
\toprule
Index sum & Monomials and their expansion weights & Total weight\\
\midrule
6 & \(ab^3/16+3a^2c^2/8\) & \(7/16\)\\
8 & \(b^4/256+3abc^2/16\) & \(49/256\)\\
9 & \(b^3c/64+ac^3/16\) & \(5/64\)\\
\bottomrule
\end{tabular}
\end{center}

Every other fiber has one monomial with weight at most one. Each fiber is thus
bounded by its largest kernel monomial. Summing proves the comparison for \(B\).

\subsubsection{The common three-point maximum for A and D}

Now put \(a=k_0\), \(b=k_1\), \(c=k_3\), and ignore \(k_2\). On the
indices \(\{0,1,3\}\), the four-fold maximum sum is exactly

\begin{equation}\label{eq:a13}
\begin{aligned}
\mathcal T(a,b,c)={}&a^4+a^3b+a^2b^2
+\max(a^3c,ab^3)+\max(a^2bc,b^4)+ab^2c\\
&+\max(a^2c^2,b^3c)+abc^2+b^2c^2+ac^3+bc^3+c^4.
\end{aligned}
\end{equation}

These terms correspond to sums \(0,1,2,3,4,5,6,7,8,9,10,12\), respectively;
sum eleven is absent. In particular \(T(k)\ge\mathcal T(a,b,c)\).

For \(A\), let
\[
\begin{aligned}
P_A={}&a^4+a^3b+a^3c+a^2b^2
+\tfrac34a^2bc+\tfrac14b^4+ab^2c\\
&+\tfrac12a^2c^2+\tfrac12b^3c
+abc^2+b^2c^2+ac^3+bc^3+c^4.
\end{aligned}
\]

Selecting one term or a convex average in each maximum in \eqref{eq:a13} gives
\(\mathcal T\ge P_A\). An exact nonnegative remainder is

\begin{equation}\label{eq:a14}
\begin{aligned}
P_A-(a+b/4+c/4)^4={}&\frac{b^2(a-b)^2}{32}
+\frac{19a^2b^2}{32}+\frac{a^2c^2}{8}\\
&+\frac{13ab^2c}{16}+\frac{13abc^2}{16}
+\frac{15ac^3}{16}+\frac{55b^4}{256}\\
&+\frac{31b^3c}{64}+\frac{125b^2c^2}{128}
+\frac{63bc^3}{64}+\frac{255c^4}{256}.
\end{aligned}
\end{equation}

Thus \eqref{eq:a2} holds for \(A\).

For \(D\), choose the terms \(ab^3,b^4,b^3c\) in the three maxima of
\eqref{eq:a13}, obtaining the lower bound

\[
P_D=a^4+a^3b+a^2b^2+ab^3+b^4+ab^2c+b^3c
+abc^2+b^2c^2+ac^3+bc^3+c^4.
\]

In comparison with \((a/4+b+c/4)^4\), the only missing monomials are

\[
\frac{a^3c}{64},\qquad\frac{3a^2bc}{16},\qquad\frac{3a^2c^2}{128}.
\]

The arithmetic--geometric mean inequalities

\[
a^3c\le\frac{3a^4+c^4}{4},\qquad
a^2bc\le\frac{a^4+b^2c^2}{2},\qquad
a^2c^2\le\frac{a^4+c^4}{2}
\]

absorb their total using at most

\[
\frac{30}{256}a^4+\frac4{256}c^4+\frac3{32}b^2c^2.
\]

The available coefficient surpluses of these three monomials in
\(P_D-(a/4+b+c/4)^4\) are respectively
\(255/256,255/256,5/8\), all larger than the required amounts. All other
coefficient surpluses are nonnegative. This proves the comparison for \(D\),
and completes \eqref{eq:a2}.

For positive \(k_i\), these comparisons are strict. The profile \(C\) has a
strict coefficient deficit at index sum two. Profiles \(A,D\) omit the positive
sum-eleven contribution involving \(k_2\); \(B\) omits the positive sum-one
contribution involving \(k_1\); and \(E\) omits the positive sum-zero
contribution. Reversal handles the other five profiles.
